\section{Method: Identifiability-Aware Source Attribution}
\label{sec:method}
\label{sec:iasa_procedure}

IASA computes from \(\At\), before fitting, diagnostics and the partition to report,
then fits once (Algorithm~\ref{alg:iasa}, with thresholds in Table~\ref{tab:thresholds}).
Contributions are estimated by nonnegative least squares (NNLS),
\begin{equation}
\widehat{\mathbf c}=\arg\min_{\mathbf c\ge 0}\;\|\yt-\At\mathbf c\|_2^2,
\label{eq:nnls_estimator}
\end{equation}
solved by projected FISTA \citep{beck2009fast} (Appendix~\ref{app:pytorch_solver}), and
it is the same estimator in every experiment.

\noindent\textbf{Diagnostics.}
IASA reports the rank of \(\At\), its smallest singular value \(\sigma_J\) (zero if rank
deficient), \(\kappa=\sigma_1/\sigma_J\), the visibility \(\|\at_j\|_2\), the weak set of
columns with visibility at most \(\tau_v\) (\(\tau_v=0\) in every run), the absorption
\(\|P_Q\mathbf a_j\|_2/\|\mathbf a_j\|_2\), which is the fraction of a column that the
projection removes, and the coherence
\(\rho_{ij}=|\at_i^\top\at_j|/(\|\at_i\|_2\|\at_j\|_2)\).  The
central diagnostic is the \emph{separation} of a group \(g\subseteq E\),
\begin{equation}
s_g=\sin\theta_g ,
\label{eq:separation}
\end{equation}
where \(\theta_g\) is the smallest principal angle \citep{bjorck1973numerical} between
\(\operatorname{col}(\At_g)\) and \(\operatorname{col}(\At_{E\setminus g})\), with
\(s_g=1\) when either space is \(\{0\}\).  In words, \(s_g\) near \(0\) means that some signal of group \(g\) is nearly reproduced
by the other sources.  For an operator with two columns \(s=\sqrt{1-\rho^2}\), and for one
column \(1/s_j^2\) is the uncentered variance-inflation factor.  IASA raises an identifiability
\emph{flag} on a merge, a weak column, an absorption at least \(0.9\),
\(\sigma_J\le10^{-6}\) or a coherence at least \(0.99\).

\noindent\textbf{Grouping.}
The partition is chosen in two steps.  First, IASA computes the finest exactly
identifiable partition \(\mathcal P^\star\) of the columns
(Theorem~\ref{thm:groupings}(c)), and the blocks of \(\mathcal P^\star\vee\mathcal D\), the
finest partition coarser than both, are the \emph{atoms}.  Second, a partition into
unions of atoms is \emph{reportable} if every block has \(s_g\ge\tau_\theta\).  IASA
returns the reportable partition with the most blocks, breaks ties by the largest
smallest separation, and lists any remaining ties.  With \(q\le12\) atoms this search is
exact by dynamic programming in \(O(3^q)\) time, whereas with more atoms it merges greedily
(Appendix~\ref{app:identifiability_proofs}).  We use \(\tau_\theta=\sqrt{1-0.99^2}=0.141\),
so for an operator with two columns the rule merges exactly when their coherence
exceeds \(0.99\).  The recorded runs used either this rule or a pairwise rule that
merges declared blocks with a cross-source coherence above \(0.99\), and on every saved
operator and observed week the two rules give the same partition.

\begin{algorithm}[t]
\caption{Identifiability-Aware Source Attribution (IASA)}
\label{alg:iasa}
\footnotesize
\begin{algsteps}
  \item \textbf{Declare} \(Q\), components, \(\mathcal D\), \(\mathcal F_0\), \(L\),
  \(\sigma_e\), \(\delta\) and thresholds.
  \item \textbf{Operator:} build \(A(\omega)\) by~\eqref{eq:transport_column}.
  \item \textbf{Project:} drop rows with a missing reading and form \(\yt\),
  \(\At\)~\eqref{eq:projected_model}.
  \item \textbf{Diagnose:} singular values, rank, visibility, absorption, coherence,
  weak set.
  \item \textbf{Group:} \(\mathcal P^\star\), then atoms \(\mathcal P^\star\vee\mathcal D\),
  then the reportable partition at \(\tau_\theta\) with the most blocks.
  \item \textbf{Fit:} solve~\eqref{eq:nnls_estimator} and form group signals
  \(\At_g\widehat{\mathbf c}_g\) and shares.
  \item \textbf{Uncertainty:} error bound \(b_g\) of each block at \(\sigma_e\), bootstrap
  intervals of the shares, and a residual adequacy test if a calibrated noise model
  exists.
  \item \textbf{Report} group signals, shares, \(b_g\), intervals, coefficients of
  single-component blocks, diagnostics and flags.
\end{algsteps}
\end{algorithm}

\noindent\textbf{Uncertainty and adequacy.}
With a declared noise standard deviation \(\sigma_e\) per reading, IASA gives for each
reported block the bound
\(b_g=\sigma_e\bigl(\sqrt{\operatorname{rank}(\At)}+\sqrt{2\log(1/\delta)}\bigr)/s_g\) on
the error of its fitted signal, which holds with probability at least \(1-\delta\)
under the conditions of Theorem~\ref{thm:group_error}(c) (\(\delta=0.05\) here).  We
call a group's signal \emph{resolved} when \(b_g\) is below the norm of its fitted
signal.  It also gives percentile intervals
of the shares from a parametric bootstrap: since the NNLS estimate depends on the
readings only through \(P_{\At}\yt\), with \(P_{\At}\) the orthogonal projector onto
\(\operatorname{col}(\At)\), the bootstrap adds Gaussian noise of standard
deviation \(\sigma_e\) in \(\operatorname{col}(\At)\) to \(\At\widehat{\mathbf c}\) and
refits.  A residual adequacy test is run when a calibrated noise model exists,
and otherwise the residual is labelled uncalibrated (Appendix~\ref{app:reporting_details}).
